\section{The covariance operator abstraction}
\label{sec:operators}

The design separates the path-following algorithm from the storage and solution of the
quadratic form (Figure~\ref{fig:design}). The turning-point loop needs only three
operations on the covariance: a matrix--vector product $\Sig x$, a multi-column solve
against the free block $\Sig_{\F\F}\inv R$, and the free-to-blocked cross product
$\Sig_{\F\B}\,x_\B$, which is formed from the first. A backend supplies them through the
protocol \code{QuadraticForm} (alias \code{CovarianceOperator}, the
\code{SymmetricOperator} contract of \pkg{cvx-linalg}), whose five members are listed in
Table~\ref{tab:protocol}; two of them serve the gross-exposure lift and the singularity
guard. The algorithm never assumes an explicit matrix, and it assumes of a backend only
that $\Sig$ is symmetric and that every free block it solves is positive definite on
the null space of the active rows, which \code{rcond\_free} lets it check before
trusting the solve (Appendix~\ref{sec:degeneracy}). Four backends ship with the package, and a user can add
others without touching the loop (\S\ref{sec:custom}).

\begin{figure}[t]
\centering
\begin{tikzpicture}[
  box/.style={draw, rounded corners=2pt, align=center, font=\small, inner sep=4pt},
  arrow/.style={-{Latex[length=2mm]}, thick}]
\node[box, minimum width=0.8\textwidth] (loop) {\textbf{Turning-point loop} (\code{CLA})\\
  first vertex $\cdot$ reduced KKT solve $\cdot$ event scan $\cdot$ least-index tie-break};
\node[box, below=0.55cm of loop, minimum width=0.8\textwidth] (proto) {\textbf{\code{QuadraticForm}}\\
  \code{n} $\cdot$ \code{matvec} $\cdot$ \code{block\_matvec} $\cdot$ \code{solve\_free} $\cdot$ \code{rcond\_free}};
\node[box, below=0.55cm of proto] (fac) {Factor\\Woodbury};
\node[box, left=0.15cm of fac] (inc) {Incremental\\maintained $\Sig_{\F\F}\inv$};
\node[box, left=0.15cm of inc] (dense) {Dense\\Cholesky};
\node[box, right=0.15cm of fac] (gram) {Gram\\$X_c\trans X_c$};
\node[box, right=0.15cm of gram, dashed] (user) {User-defined\\(\S\ref{sec:custom})};
\draw[arrow] (loop) -- (proto);
\foreach \b in {dense, inc, fac, gram, user} \draw[arrow] (proto.south) -- (\b.north);
\end{tikzpicture}
\caption{The covariance enters the turning-point loop only through the
\code{QuadraticForm} protocol; the backends below it are interchangeable.}
\label{fig:design}
\end{figure}

\begin{table}[t]
\centering
\small
\begin{tabular}{p{0.32\textwidth}p{0.29\textwidth}p{0.28\textwidth}}
\hline
Member & Returns & Used for \\
\hline
\code{n} & the number of assets & sizes \\
\code{matvec(x)} & $\Sig x$, for a vector or a matrix $x$ & gradients, the cross product \\
\code{block\_matvec(rows, cols, v)} & $\Sig_{\text{rows},\text{cols}}\, v$ & the gross-exposure lift (\S\ref{sec:leverage}) \\
\code{solve\_free(free, R)} & $\Sig_{\F\F}\inv R$ & the reduced KKT solve (\S\ref{sec:blockelim}) \\
\code{rcond\_free(free)} & the reciprocal condition number of $\Sig_{\F\F}$ & the singularity guard \\
\hline
\end{tabular}
\caption{The \code{QuadraticForm} protocol. Index sets are integer arrays into $\{0,\dots,n-1\}$.}
\label{tab:protocol}
\end{table}

\paragraph{Dense.} \code{DenseCovariance} wraps an explicit matrix and solves the free
block by Cholesky (\code{scipy.linalg.cho\_factor}), falling back to LU on the rare
block that is not numerically positive definite. It reproduces the classical CLA and is
the reference backend.

\paragraph{Incremental dense.}
\label{sec:incremental}
\code{IncrementalDenseCovariance} traces the same turning points but \emph{maintains}
$\Sig_{\F\F}\inv$: because the free set changes by one asset per turning point, a
bordered rank-one update or a symmetric deletion costs $O(n_\F^2)$ against a fresh
$O(n_\F^3)$ factorisation, which makes a trace about twice as fast at $n=320$, three times at $n=640$, and seven times at $n=2560$ (\S\ref{sec:baseline}). A non-positive Schur pivot or a multi-asset change triggers a full refactorisation. The
backend passes the same KKT certificate as the dense one (Table~\ref{tab:kkt}), but a
maintained inverse carries round-off from step to step over the roughly $n$ updates of
a trace, so it is opt-in.

\paragraph{Factor (diagonal-plus-low-rank).}
\code{FactorCovariance} represents
\begin{equation}
\label{eq:factor}
\Sig = \operatorname{diag}(d) + U\,\Delta\,U\trans,
\qquad d\in\R^{n}_{>0},\; U\in\R^{n\times K},\; \Delta\in\R^{K\times K},
\end{equation}
the form of factor models \citep{sharpe1963} and of eigenvalue-clipped covariances. The
free-block solve never forms $\Sig_{\F\F}$; it uses the Woodbury identity
\begin{equation}
\label{eq:woodbury}
\Sig_{\F\F}\inv = D_\F\inv - D_\F\inv U_\F\, W\inv\, U_\F\trans D_\F\inv,
\qquad
W = \Delta\inv + U_\F\trans D_\F\inv U_\F \in\R^{K\times K},
\end{equation}
so a solve costs $O(n_\F K^2 + K^3)$ rather than $O(n_\F^3)$, storage is $O(nK)$, and
the cross product is $U_\F\,\Delta\,(U_\B\trans x_\B)$. The same free-block algebra serves a
single constrained solve as well: \citet{schmelzer2026nncg} reach the free block of a
non-negative QP through the same Woodbury inverse and factored products.

The admissible model is $d > 0$ and $\Delta$ symmetric positive semidefinite.
\code{FactorCovariance} folds $\Delta$ into the loadings when it is built: with
$\Delta = V\Lambda V\trans$ it stores $U' = U V_+ \Lambda_+^{1/2}$, the directions of
positive variance, and an identity inner matrix, which is the same $\Sig$. The
Woodbury solve \eqref{eq:woodbury} then runs with $U'$ and $W = I + U_\F'^{\mathsf{T}}
D_\F\inv U_\F'$ and never inverts $\Delta$. A singular factor covariance, from a factor
with zero variance or perfectly correlated factors, is therefore admissible, and its rank
$K$ is the number of directions kept. An indefinite $\Delta$ is refused: a signed
low-rank term can still give a positive-definite $\Sig$, but $W$ would then be
indefinite, so such a covariance is passed as a dense matrix. Since $d > 0$, $W$ is
positive definite, and the factor solve needs no guard of its own beyond the one on
$\Sig_{\F\F}$. Against the dense trace of the same $\Sig$, the factor trace has the same
turning points and weights equal to $2.8\times10^{-12}$, with KKT residuals below
$2.3\times10^{-13}$, for $\Delta$ with condition numbers up to $10^{14}$, loadings
collinear to within $10^{-10}$, and the full rank $K = n = 120$
(\target{validate-factor}). Every active constraint row adds a column to the solve and a
row to the Schur complement, so with many active rows that term dominates the cost. A low-rank \emph{approximation} of a sample covariance traces the frontier of
the approximation, which is a different frontier (\S\ref{sec:estimators}).

\paragraph{Data matrix (Gram).}
\label{sec:gram}
A sample covariance is a Gram matrix: for the column-centred returns
$X_c\in\R^{T\times n}$, $\Sig = X_c\trans X_c/(T-1)$. \code{GramCovariance} implements
the protocol straight from $X_c$ in $O(Tn)$ storage. When $T<n$ the matrix is singular,
and once the free set outgrows the data rank the free block is singular too; the active
rows still hold it on their null space, where it is solved (Appendix~\ref{sec:degeneracy}).
An optional ridge $\theta I$ (the argument \code{ridge}) makes $\Sig$ positive definite,
a different, regularised frontier, and the solve is then Woodbury in the
$T$-dimensional observation space. The saving is in storage: when the dense matrix fits,
\code{DenseCovariance} is faster per step, and the structured route pays only while $T$
(or $K$) stays well below $n$ (\S\ref{sec:rankscaling}).

\subsection{Regularization through the operator}
\label{sec:regularization}

A quadratic penalty $\tfrac{\theta}{2}\, w\trans M w$, with $\theta \ge 0$ and $M$
symmetric positive semidefinite and independent of $\lambda$, changes only the Hessian,
so the path stays piecewise linear and the penalty is the substitution
$\Sig \mapsto \Sig + \theta M$: a ridge ($M = I$) or a factor-exposure penalty
($M = B B\trans$ for exposures $B$). \code{CLA} has no penalty argument; the regularized
form is supplied as the covariance, a bundled or a user-written \code{QuadraticForm}
(\S\ref{sec:custom}). \code{GramCovariance} with $\theta>0$ and \code{FactorCovariance}
already are such forms. A rank-deficient sample covariance needs neither: its frontier
is traced as it is, as long as the active rows keep the free block positive definite on
their null space (Appendix~\ref{sec:degeneracy}). A small ridge or a factor model is the
fix when they do not, and a choice of estimator otherwise (\S\ref{sec:estimators}). A penalty centred on a nonzero portfolio $w_0$ (tracking
error, quadratic turnover) adds the linear term $-\theta\, w\trans M w_0$, which the
interface does not expose; nonquadratic penalties such as linear transaction costs lie
outside the method.

\subsection{Writing a backend}
\label{sec:custom}

A backend is a subclass of \code{QuadraticForm} that implements the five members of
Table~\ref{tab:protocol}. For a block-diagonal covariance, assets in groups with no
cross-group covariance (a structure the package does not ship), the free-block solve
splits into one Cholesky solve per block. The two members that carry the structure are
below; \code{n}, \code{matvec} and \code{block\_matvec} are a few lines each, and the
complete, runnable class is in the package documentation
(\srcfile{docs/custom-backend.md}{docs/custom-backend.md}).

\begin{CodeInput}
class BlockDiagonal(QuadraticForm):   # Sigma = blockdiag(S_1, ..., S_m)
    def solve_free(self, free, rhs):
        out = np.empty_like(rhs, dtype=float)
        # one (positions, free sub-block) pair per group
        for pos, sub in self._parts(free):
            out[pos] = cho_solve(cho_factor(sub), rhs[pos])
        return out

    def rcond_free(self, free):
        eig = np.concatenate(
            [np.linalg.eigvalsh(sub) for _, sub in self._parts(free)]
        )
        return eig.min() / eig.max()
\end{CodeInput}

\code{CLA} runs on it unchanged: on four groups of five assets it traces the same $19$
turning points as the dense matrix. The package's test suite,
\srcfile{tests/test_custom_backend.py}{tests/test\_custom\_backend.py}, checks the same
with general constraints and with a gross-exposure cap, which goes through
\code{block\_matvec}.
The frontier is exact to the accuracy of \code{solve\_free}; an approximate solver
traces the frontier of the matrix it effectively applies, and the KKT certificate of
\S\ref{sec:exactness} checks the result. A sparse matrix
with a sparse Cholesky factor, a Kronecker product, or a matrix-free operator with an
iterative \code{solve\_free} fit the same five members.
